\section{Functional Gaussian limits and the physical clock}
\label{sec:functional-gaussian}
A central-frequency estimate alone does not give tightness of the record
process. We prove it by smoothing on a slower scale, using fourth moments
for the smooth sums and a maximal second-moment estimate for the residual.
The estimates remain uniform over the moving physical family.

\subsection{Fourth moments and removal of smoothing}
\begin{lemma}[A fourth-moment bound at a specified smoothing scale]
\label{lem:smooth-fourth-moment}
For $m\ge1$, $R\in I$, and $0<\delta<1/2$,
\begin{equation}\label{eq:smooth-fourth-moment}
 \int |S_{m,R}h_{R,\delta}|^4\dd\nu
       \le C\bigl(m^2+(m+1)\delta^{-8}\bigr).
\end{equation}
For $r_{R,\delta}=h_R-h_{R,\delta}$ and $L\ge1$,
\begin{equation}\label{eq:residual-process-maximal}
 \int\max_{0\le j\le L}|S_{j,R}r_{R,\delta}|^2\dd\nu
 \le CL\delta(1+|\log\delta|)[\log(2+L)]^2.
\end{equation}
Both bounds hold, with changed constants, after multiplication by any
nonnegative initial density bounded by a fixed constant.
\end{lemma}
\begin{proof}
For the first assertion consider one real component $u$ of
$h_{R,\delta}$ and the one-variable family in
Lemma~\ref{lem:smooth-collision-expansion}. Write
$\psi(z)=\log\lambda(z)$ and $A(z)=\ell_j(\Pi(z)\nu)$. The spectral
pairing for the stationary characteristic function is
\[
 \int e^{izS_{m,R}u}\dd\nu=e^{m\psi(z)}A(z)+E_m(z).
\]
On a complex disk of radius $a\delta^2$, with $a$ decreased once,
$|E_m(z)|\le C\rho^m$, $|A(z)|\le C$, and $|\psi(z)|\le C$.
Cauchy's inequalities therefore give
$|A^{(j)}(0)|\le C_j\delta^{-2j}$ for $j\le4$,
$|\psi^{(j)}(0)|\le C_j\delta^{-2j}$ for $j=3,4$, and
$|E_m^{(4)}(0)|\le C\delta^{-8}\rho^m$.
In addition $A(0)=1$, $\psi'(0)=0$, and
$|\psi''(0)|\le C$, by the uniformly summable smoothed covariance.
The fourth derivative of the principal term at zero is exactly
\[
 A^{(4)}(0)+6m\psi''(0)A''(0)+4m\psi'''(0)A'(0)
       +m\psi^{(4)}(0)+3m^2\psi''(0)^2.
\]
It is bounded by $C(m^2+(m+1)\delta^{-8})$. Differentiating under the
finite-time integral is legitimate because $u$ is bounded. The fourth
derivative is its fourth moment. Summing the component estimates proves
\eqref{eq:smooth-fourth-moment}.

For the second assertion, \eqref{eq:small-bv-variance} bounds the second
moment of every interval sum of length $\ell$ by
$C\ell\delta(1+|\log\delta|)$. Apply the dyadic proof of
Lemma~\ref{lem:dyadic-collision-maximal} with this constant. The final
assertion follows by domination of the initial density; it does not
assert stationarity under that density.
\end{proof}

For $A_0<\infty$, let $X_{n,R}$ be the polygonal interpolation on
$[0,A_0]$ with values
\[
 X_{n,R}(j/n)=n^{-1/2}S_{j,R}h_R.
\]
Let $\mathcal W_C$ denote Brownian motion starting at zero with covariance
$\mathbb E[\mathcal W_C(s)\mathcal W_C(t)^{\mathsf T}]=(s\wedge t)C$;
$C$ is allowed to be singular. We use the bounded-Lipschitz distance on
probability laws on continuous path space, with its supremum norm.

\begin{theorem}[Uniform collision functional central limit]
\label{thm:collision-functional}
Fix $A_0,B_0<\infty$. Uniformly for $R\in I$ and all nonnegative
probability densities $\sigma_R$ satisfying
$\|\sigma_R\|_\infty+\|\sigma_R\|_{\mathrm{BV}}\le B_0$,
\begin{equation}\label{eq:collision-functional}
 d_{\rm BL}\bigl(\operatorname{Law}_{\sigma_R\nu}(X_{n,R}),
                     \operatorname{Law}(\mathcal W_{\Gamma_R})\bigr)
       \longrightarrow0.
\end{equation}
In particular this holds under $\nu$ and under $\nu_R^*$.
\end{theorem}
\begin{proof}
Set $\delta_n=\tfrac14 n^{-1/32}$ and let $X_{n,R}^{\delta_n}$ be the
polygonal process with $h_R$ replaced by $h_{R,\delta_n}$.
Lemma~\ref{lem:smooth-fourth-moment} implies
\[
 \sup_{R,\sigma_R}\int\|X_{n,R}-X_{n,R}^{\delta_n}\|_\infty^2
             \sigma_R\dd\nu
       \le C_{A_0,B_0}\delta_n[\log(2+n)]^3\longrightarrow0.
\]
Integer parts and interpolation change only the constant. Thus it
suffices to establish the limit for the smoothed process. This scale
is slower than the one optimizing the one-time estimate.

We first prove finite-dimensional convergence. Fix
$0=t_0<t_1<\cdots<t_q\le A_0$ and vectors $v_1,\ldots,v_q$.
Put $m_j=\lfloor nt_j\rfloor-\lfloor nt_{j-1}\rfloor$ and
$z_j=v_j/\sqrt n$. Smooth the initial density at the same scale;
its $L^1$ error is $O(\delta_n)$ and its distribution norm is
$O(\delta_n^{-2})$. The characteristic function of the consecutive
increments, before the harmless polygonal endpoint corrections, is
\[
 \ell_j\bigl(\mathcal L_{R,\delta_n,z_q}^{m_q}\cdots
       \mathcal L_{R,\delta_n,z_1}^{m_1}
                   (\mathsf S_{\delta_n}\sigma_R)\nu\bigr).
\]
The product is chronological, with the earliest segment on the right.
For each fixed $q$, substitute the spectral decompositions of
Lemma~\ref{lem:smooth-collision-expansion}. Every term with a
complementary power tends to zero uniformly, since $m_j$ is a positive
multiple of $n$ and its remaining factors are polynomially bounded in
$\delta_n^{-1}$. The product of principal amplitudes tends to one:
$\Pi(0)=\nu\ell_j$, and its error is
$O_q(\delta_n^{-4}n^{-1/2})$. The cubic error in the logarithm is
$O_q(\delta_n^{-6}n^{-1/2})=o(1)$. The resulting characteristic function
is therefore, uniformly in $R$ and in the initial densities,
\[
 \exp\left(-\frac12\sum_{j=1}^q(t_j-t_{j-1})
            v_j^{\mathsf T}\Gamma_Rv_j\right)+o(1).
\]
Here \eqref{eq:collision-covariance-limit} removes smoothing from the
covariance. This is the joint law of the independent Brownian
increments, not merely a collection of marginal Gaussian laws.

For tightness put $L_n=\lceil\delta_n^{-8}\rceil$ and interpolate the
smoothed sums only at collision indices $0,L_n,2L_n,\ldots$.
Call this coarse polygonal process $Z_{n,R}$. An interval of $m\ge L_n$
collisions has fourth moment at most $Cm^2$, by
\eqref{eq:smooth-fourth-moment} and stationarity of $\nu$. Domination
by $\sigma_R\le B_0$ gives the same bound from every deterministic
starting index under the initial law. It follows at coarse grid points
that the fourth moment of an increment is at most $C|t-s|^2$.
For arbitrary points, split an increment into its two fractional end
segments and the complete grid interval. Within one grid interval,
linear interpolation gives a factor $(|t-s|n/L_n)^4$ multiplying a
fourth moment $O((L_n/n)^2)$, which is bounded by $C|t-s|^2$.
The same estimate follows across adjacent intervals and then across
any number of intervals, with a fixed constant. Hence
\[
 \mathbb E_{\sigma_R\nu}|Z_{n,R}(t)-Z_{n,R}(s)|^4
       \le C_{A_0,B_0}|t-s|^2.
\]
The usual dyadic proof of the continuity and tightness criterion applies
uniformly to these continuous processes. Extend the last grid interval
past $A_0$ and restrict back to $[0,A_0]$ to avoid an endpoint convention.
Finally, boundedness of $h_{R,\delta_n}$ gives deterministically
\[
 \|Z_{n,R}-X_{n,R}^{\delta_n}\|_\infty
       \le C L_n/\sqrt n=O(n^{-1/4})\longrightarrow0.
\]
This proves tightness of the original processes as well.

To obtain the stated uniform distance, suppose it failed and select a
violating sequence of radii and initial densities. By compactness a
subsequence of the radii converges to $R_0$. The uniform tightness and
finite-dimensional limits identify every subsequential path law as
$\mathcal W_{\Gamma_{R_0}}$. Continuity of $\Gamma_R$, including at
singular matrices, gives continuity of the Brownian law; for example
couple by $\Gamma_R^{1/2}$ times one standard Brownian motion. This
contradicts the violating bounded-Lipschitz distance and proves
\eqref{eq:collision-functional}.
\end{proof}

\subsection{The induced record process}
Let $Y_{n,R}$ be the polygonal interpolation on $[0,1]$ of
$n^{-1/2}(J_{j,R}-j\bar G_R)$ at times $j/n$.

\begin{theorem}[Functional limit of the actual return process]
\label{thm:actual-return-functional}
Uniformly for $R\in I$,
\begin{equation}\label{eq:actual-return-functional}
 d_{\rm BL}\bigl(\operatorname{Law}_{\nu_R^*}(Y_{n,R}),
                         \operatorname{Law}(\mathcal W_{D_R})\bigr)
       \longrightarrow0,
       \qquad D_R=c_*^{-1}\Gamma_R.
\end{equation}
\end{theorem}
\begin{proof}
The count coordinate of Theorem~\ref{thm:uniform-return-fluid} gives
\[
 \sup_{0\le s\le1}|N_{\lfloor ns\rfloor,R}/n-s/c_*|
                  \longrightarrow0
\]
in probability, uniformly in $R$. Choose $A_0>c_*^{-1}+1$ and apply
Theorem~\ref{thm:collision-functional} on $[0,A_0]$ under $\nu_R^*$.
Tightness with continuous limits gives a uniform probabilistic modulus
of continuity for $X_{n,R}$. Consequently replacing its random argument
$N_{\lfloor ns\rfloor,R}/n$ by $s/c_*$ changes it by $o_P(1)$ in
supremum norm. The random argument remains in $[0,A_0]$ with probability
tending to one uniformly; independence from $X_{n,R}$ is not required.
The exact identity \eqref{eq:exact-stopped-compensation} identifies this
composition with the step version of the centered induced process.

For its polygonal version, stationarity under $F_R^*$ and the true
one-block tail give, for every $\epsilon>0$,
\[
 \nu_R^*\left\{\max_{j\le n}r_R^*\circ(F_R^*)^j>
                  \epsilon\sqrt n\right\}
       \le C(n+1)e^{-c\epsilon\sqrt n}\longrightarrow0.
\]
Each centered return increment has norm at most $C r_R^*$, by
boundedness of $h_R$ and exact compensation. Thus interpolation changes
the process by $o_P(1)$ uniformly. Brownian time change $s\mapsto s/c_*$
has covariance $c_*^{-1}\Gamma_R$, which proves the assertion by the
same compactness argument as above.
\end{proof}

\subsection{An unconditioned physical-time consequence}
Recall the actual stationary physical collision count $C_R(t)$ and
lattice displacement $W_R(t)$ of
Theorem~\ref{thm:uniform-physical-clock}. Define the linear map
\[
 B_R(x_1,x_2,x_3,x_4)=(x_1,x_2,x_3-x_4/\bar\tau_R),\qquad
 \mathcal V_R=\bar\tau_R^{-1}B_R\Gamma_R B_R^{\mathsf T}.
\]
\begin{corollary}[Physical displacement and collision fluctuations]
\label{cor:physical-functional-gaussian}
Under the stationary laws $\mu_R$, the processes
\begin{equation}\label{eq:physical-functional-gaussian}
 t^{-1/2}\bigl(W_R(ts),C_R(ts)-ts/\bar\tau_R\bigr),\quad 0\le s\le1,
\end{equation}
converge uniformly in $R$ to Brownian motion with covariance $\mathcal V_R$.
The convergence may be taken in the Skorokhod space, or for continuous
interpolations whose uniform difference from the displayed processes
is $O(t^{-1/2})$. The matrices $\mathcal V_R$ are continuous and positive
semidefinite. No conditioning event is changed or imposed in this result.
\end{corollary}
\begin{proof}
In the physical collision suspension, the marginal of the last outgoing
collision state $x$ is $\sigma_R\nu$, where
$\sigma_R=\tau_R/\bar\tau_R$. Its supremum and variation norms are
uniformly bounded by Lemma~\ref{lem:physical-bv} and
$\min_I\bar\tau_R>0$. The elapsed time since that collision is bounded
by the uniform single-flight horizon. Thus
Theorem~\ref{thm:collision-functional} applies with this initial density.

The section compensation disappears under $B_R$:
\[
 B_Rh_R=(\kappa_{{\rm coll},1},\kappa_{{\rm coll},2},
                    1-\tau_R/\bar\tau_R),
 \qquad B_R\bar G_R=0.
\]
The collision clock satisfies
$\sup_{s\le1}|C_R(ts)/t-s/\bar\tau_R|=o_P(1)$ uniformly, by
Theorem~\ref{thm:uniform-physical-clock}. Apply the continuous-limit
random-time argument to $B_RX_{\lfloor t\rfloor,R}$ on a fixed interval
larger than $1/\min_I\bar\tau_R+1$. Rounding $t$ affects neither the
normalization nor the limit.

The initial and final partial physical flights have uniformly bounded
durations. The collision sum of the lattice labels differs from
$W_R(ts)$ by a uniformly bounded cell-position term. Also the sum of
completed flight times differs from $ts$ by a uniformly bounded amount.
Consequently the stopped sum of $B_Rh_R$ differs from the numerator in
\eqref{eq:physical-functional-gaussian} by a uniformly bounded vector,
simultaneously for all $s$. The jumps of the lattice and collision counts
are uniformly bounded; interpolation between collision times, together
with the same cell-position bound, has error $O(t^{-1/2})$. The limiting
Brownian time change is $s/\bar\tau_R$, giving $\mathcal V_R$. Continuity follows
from that of $\Gamma_R$ and the explicit strictly positive mean roof.
\end{proof}

The functional results supply the unconditioned process input for the
original clock analysis. They neither evaluate a measurable transfer
function at the selected periods nor control a shrinking conditioning
event. Uniform positive definiteness is supplied separately by
Corollary~\ref{cor:actual-uniform-ellipticity}; the weighted raw density
estimates remain distinct from these path-space conclusions.
